\documentclass[10pt,a4paper]{amsart}
\usepackage[margin=19mm]{geometry}
\usepackage{amsmath,amssymb,mathtools}
\usepackage[scr=rsfs,cal=euler]{mathalfa}
\usepackage{xfrac}
\usepackage[unicode,hidelinks,bookmarks=false]{hyperref}
\newcommand{\ie}{i.e.\ }
\newcommand{\cf}{cf.\ }
\newcommand{\resp}{respectively\ }
\newcommand{\Subsection}{\subsection}
\providecommand{\nobreakdash}{\nobreak\hbox{-}}
\setlength{\emergencystretch}{2em}
\multlinegap0pt
\usepackage{amssymb}
\usepackage[shortlabels]{enumitem}
\usepackage{mathtools}
\usepackage{bbm}
\usepackage{bm}
\newtheorem{theorem}{Theorem}
\newtheorem{corollary}{Corollary}
\newcommand{\R}{\mathbb{R}}
\newcommand{\wh}{\widehat}
\newcommand{\wt}{\widetilde}
\title{Minimal Solutions to the Skorokhod Reflection Problem Driven by Jump Processes and an Application to Reinsurance}
\author{Graeme Baker, Ankita Chatterjee}
\date{Source: arXiv:2512.24491v1 (CC BY 4.0)}
\begin{document}
\maketitle

\noindent\emph{Source and licence.} Minimal Solutions to the Skorokhod Reflection Problem Driven by Jump Processes and an Application to Reinsurance. Version arXiv:2512.24491v1 (CC BY 4.0), distributed by arXiv under the Creative Commons Attribution 4.0 International licence, http://creativecommons.org/licenses/by/4.0/, as recorded in the arXiv licence metadata for this exact version and shown on the abstract page https://arxiv.org/abs/2512.24491v1. Full licence notice: \texttt{licenses/arxiv-2512.24491-CC-BY-4.0.txt}.\par\smallskip

\noindent\textit{Editorial scope.} Counted unit: the article's corollary coro (source0.tex lines 146--148). The article's complete original proof of it is delivered with it (source0.tex lines 155--182, one \texttt{proof} environment of 28 lines, which the article places away from the statement). No mathematics is written, repaired or re-derived: the statement and the proof are the article's own lines, in the article's own notation, and no retained line is substituted or re-flowed.  Retained uncounted as the source-local context the proof needs: lines 114--116; lines 122--124; lines 139--142.  Nothing was omitted from any retained span, so the package carries no omission record.  No retained line contains a citation, so the wrapper carries no bibliography and no bibliography entry.  No retained line contains a figure environment, an image-inclusion command or drawing code, so no image is delivered and the package declares no asset dependency.  Editorial formatting only, in addition: an AMS \texttt{amsart} preamble that restates the article's own declarations and macros used by the retained lines, namely the environments \texttt{theorem}, \texttt{corollary} on separate counters, together with the standard packages those lines need.  The retained environments print in this document as Theorem 1, Corollary 1 in order of appearance, whereas the article prints the same items with the numbering of its own full document.  The complete native source of the article as distributed in the arXiv e-print for this exact version is bundled unchanged as \texttt{sources/191902/source0.tex}.
\begin{equation}\label{eq:ps}
X_t^i=X_{0-}^i+c_it-Z_t^i+L_t^i-\sum_{j\neq i}q_{ij}L_t^j,\quad 1\le i\le n,
\end{equation}

\begin{equation}\label{eq:Ldef}
L_t^i=\sup_{s\le t}\bigg(X_{0-}^i+c_it-Z_s^i-\sum_{j\neq i}q_{ij}L_s^j\bigg)_-,\quad 1\le i\le n,
\end{equation}

\begin{theorem}\label{thm}
Consider the cone $C=\{u\in\R^n:u\ge 0, u^\top (I-Q)\le 0\}$ and its dual cone $C^*=\{y\in\R^n:u^\top y \ge 0 \text{ for all }u\in C\}$.
A solution $(X,L)$ of \eqref{eq:ps}--\eqref{eq:Ldef} on an interval $[0,\tau)$ can be extended to $[0,\tau]$ if and only if $X_{\tau-}-\Delta Z_\tau \in C^*$. Moreover, when $X_{\tau-}-\Delta Z_\tau \in C^*$ there exists a unique minimal $\Delta L_\tau$ with respect to the partial order $\ge$ on $\R_+^n$ such that $L_\tau=L_{\tau-}+\Delta L_\tau$ gives an extension of $(X,L)$ on $[0,\tau]$.
\end{theorem}

\begin{corollary}\label{coro}
Given inputs $X_{0-}^1,\dots, X_{0-}^n$, and $Z^1,\dots,Z^n$, there exists unique $(X,L)$ and $\tau^*$ such that $(X,L)$ solves \eqref{eq:ps}--\eqref{eq:Ldef} on $[0,\tau^*)$ and if $(\wt{X},\wt{L})$ is any other strong solution to \eqref{eq:ps}--\eqref{eq:Ldef} on $[0,\wt\tau)$ with the same inputs as $(X,L)$ then $\wt\tau\le\tau^*$ and $\wt{L_t}\ge L_t$ for all $t\in[0,\wt\tau)$.
\end{corollary}

\begin{proof}[Proof of Corollary \ref{coro}]
Consider any arbitrary solution  $(\wt{X},\wt{L})$ to \eqref{eq:ps}--\eqref{eq:Ldef} on some interval $[0,\wt\tau)$.
Let $\tau_1$ denote the first time that one or multiple of the processes $Z^1,\dots,Z^n$ exhibit a jump. Set $L^1_t=\dots =L^n_t=0$ on $[0,\tau_1)$. From \eqref{eq:Ldef}, we see that for $1\le i\le n$, $\wt{L}^i_t\ge 0=L^i_t$ for all $t\in [0,\tau_1\wedge \wt\tau)$, where we have used the notation $s\wedge t:=\min(s,t)$. We consider now two cases:

\medskip

\emph{Case 1.} If $X_{\tau_1-}-\Delta Z_{\tau_1} \notin C^*$, then we set $\tau^*=\tau_1$. If $\wt\tau \ge \tau^*$ then we must have $\wt{X}_{\tau_1-}-\Delta Z_{\tau_1} \notin C^*$ and hence $\wt\tau > \tau^*$ is impossible.

\medskip

\emph{Case 2.} If $X_{\tau_1-}-\Delta Z_{\tau_1} \in C^*$ we let $\Delta L_{\tau_1}$ be the unique minimal jump size from Theorem \ref{thm} and set $L_{\tau_1}=0+\Delta L_{\tau_1}$.
If $\wt\tau < \tau_1$, then $\wt{L}^i_t\ge L^i_t$ for all $t\in [0,\tau_1\wedge \wt\tau]$ trivially. If not, then suppose for the sake of contradiction that $\wt{L}^i_{\tau_1}< L^i_{\tau_1}$ for some $1\le i\le n$ and for $1\le i \le n$ set 
\begin{equation*}
\wh{L}^i_t=\begin{cases}
0&t<\tau_1\\
\wt{L}^i_{\tau_1} & t=\tau_1.
\end{cases}
\end{equation*}
Then, the definition of $\tau_1$ along with the non-negativity of $\wt{X}$ imply that
\begin{equation*}
X_{0-}^i+c_it-Z^i_t+\wh{L}^i_t-\sum_{j\neq i}q_{ij} \wh{L}^j_t\ge 0
\end{equation*}
for all $t\in[0,\tau_1]$ and $1\le i\le n$. However, $(\Delta\wh{L}^1_{\tau_1},\dots,\Delta\wh{L}^n_{\tau_1})=(\wt{L}^1_{\tau_1},\dots, \wt{L}^n_{\tau_1})$ contradicts the minimality of $\Delta L_{\tau_1}$ with respect to the partial order $\ge$ on $\R_+^n$. Therefore $\wt{L}^i_t\ge L^i_t$ for all $t\in [0,\tau_1\wedge \wt\tau]$. 

\medskip

Let $\tau_2$ denote the next jump time for any of the processes $Z^1,\dots,Z^n$. Repeating the above argument shows that $\wt{L}^i_t\ge L^i_t$ for all $t\in [0,\tau_2\wedge \wt\tau]$ and $\wt\tau > \tau^*$ is impossible. We continue as before with $\tau_3,\tau_4$, and so on. Since $(\wt{X},\wt{L})$ and $\wt\tau$ are arbitrary, we have completed the proof. \qedhere
\end{proof}

\end{document}
